\section{Introducción}
\label{cap:resultados}

Se reportan los resultados de la línea base (\eegnet{}) y de las tres arquitecturas
objetivo: el Conformer, el modelo fundacional LaBraM (variante Base) y EEGMamba. Con la media
de tres semillas, el protocolo \loso{} supera al intra-sujeto en las cuatro arquitecturas,
\eegnet{}, el Conformer y LaBraM son equivalentes entre sí dentro de $\pm 0.03$ y solo LaBraM
muestra una degradación significativa ante participantes no vistos. Después se informan tres
grupos de variantes sobre la misma base: preprocesamiento y aumentación en el protocolo
intra-sujeto (subsección~\ref{sec:ablaciones-f4}), alineamiento euclidiano y reflexión de canales
en \loso{} (subsección~\ref{sec:f6f8}) y calibración con pocos ensayos del sujeto no visto
(subsección~\ref{sec:fewshot}). Sus condiciones de lectura están en la
subsección~\ref{sec:condiciones-variantes}.



Las cuatro pruebas examinaron la generalización entre personas (Prueba 1), las variantes de procesamiento y calibración (Prueba 2), la transferencia entre conjuntos (Prueba 3) y los criterios adicionales al desempeño (Prueba 4). En conjunto aportaron evidencia para el objetivo~\ref{oe:4}; la segunda examinó el procesamiento del objetivo~\ref{oe:2}, y la primera y la cuarta evaluaron los modelos del objetivo~\ref{oe:3}.
